\documentclass[12pt]{article}

\usepackage{fullpage}
\usepackage{multicol,multirow}
\usepackage{tabularx}
\usepackage{ulem}
\usepackage{hyperref}
\usepackage{graphicx}
\usepackage[utf8]{inputenc}
\usepackage[russian]{babel}

\begin{document}

\section*{Лабораторная работа №\,2 по курсу дискрeтного анализа: сбалансированные деревья}

Выполнил студент группы М8О-207БВ-24 МАИ \textit{Устинов Александр}.

\subsection*{Условие}
 
\begin{enumerate}
\item Общая постановка задачи.
Реализовать декартово дерево с возможностью поиска, добавления и удаления элементов.

Необходимо создать программную библиотеку, реализующую указанную структуру данных, на основе которой разработать программу-словарь. В словаре каждому ключу, представляющему из себя регистронезависимую последовательность букв английского алфавита длиной не более 256 символов, поставлен в соответствие некоторый номер. 
Разным словам может быть поставлен в соответствие один и тот же номер.

Программа должна обрабатывать строки входного файла до его окончания. Каждая строка может иметь следующий формат:

+ word 34 — добавить слово «word» с номером 34 в словарь. Программа должна вывести строку «OK», если операция прошла успешно, «Exist», если слово уже находится в словаре.

- word — удалить слово «word» из словаря. Программа должна вывести «OK», если слово существовало и было удалено, «NoSuchWord», если слово в словаре не было найдено.

word — найти в словаре слово «word». Программа должна вывести «OK: 34», если слово было найдено; число, которое следует за «OK:» — номер, присвоенный слову при добавлении. В случае, если слово в словаре не было обнаружено, нужно вывести строку «NoSuchWord».

! Save /path/to/file — сохранить словарь в бинарном компактном представлении на диск в файл, указанный парамером команды. В случае успеха, программа должна вывести «OK», в случае неудачи выполнения операции, программа должна вывести описание ошибки (см. ниже).

! Load /path/to/file — загрузить словарь из файла. Предполагается, что файл был ранее подготовлен при помощи команды Save. В случае успеха, программа должна вывести строку «OK», а загруженный словарь должен заменить текущий (с которым происходит работа); в случае неуспеха, должна быть выведена диагностика, а рабочий словарь должен остаться без изменений. Кроме системных ошибок, программа должна корректно обрабатывать случаи несовпадения формата указанного файла и представления данных словаря во внешнем файле.

Для всех операций, в случае возникновения системной ошибки (нехватка памяти, отсутсвие прав записи и т.п.), программа должна вывести строку, начинающуюся с «ERROR:» и описывающую на английском языке возникшую ошибку.

\item Вариант задания.
Структура данных: AVL-дерево.
\end{enumerate}

\subsection*{Метод решения}

Для решения задачи была реализована структура данных <<AVL-дерево>> — самобалансирующееся бинарное дерево поиска, которое поддерживает логарифмическую сложность операций вставки, удаления и поиска.
Балансировка осуществляется путём выполнения поворотов при нарушении баланса узла, что гарантирует высоту дерева $O(\log n)$.

\begin{enumerate}
\item Вставка нового узла: рекурсивный спуск по дереву, затем балансировка с проверкой коэффициента баланса.
\item Удаление узла: поиск удаляемого элемента, замена на минимальный из правого поддерева (при необходимости), балансировка.
\item Поиск узла: рекурсивный спуск по сравнению ключей до нахождения совпадения или достижения листа.
\item Балансировка: четыре случая поворотов (LL, LR, RR, RL) в зависимости от знака баланса узла и его потомков.
\end{enumerate}

\textbf{Архитектура программы:}

\begin{center}
\begin{tabular}{@{}l@{}}
\texttt{main.cpp} \\
\quad $\vdash$ \textbf{Ввод данных} --- чтение из stdin через \texttt{cin} (команды: \texttt{+}, \texttt{-}, \texttt{!}, поиск) \\
\quad $\vdash$ \textbf{Структура \texttt{Node}} --- узел AVL-дерева: ключ (char*), значение (uint64\_t), высота, потомки \\
\quad $\vdash$ \textbf{Класс \texttt{Dictionary}} --- AVL-дерево с методами вставки, удаления, поиска, сохранения/загрузки \\
\quad $\vdash$ \textbf{Балансировка} --- повороты (левый/правый) при нарушении баланса узла \\
\quad $\vdash$ \textbf{Вывод} --- запись в stdout через \texttt{cout} (OK, Exist, NoSuchWord, ERROR) \\
\end{tabular}
\end{center}

\vspace{0.3cm}


\textbf{Использованные источники:}
\begin{itemize}
\item Кормен Т., Лейзерсон Ч., Ривест Р., Штайн К. ``Алгоритмы: построение и анализ''.
\item \url{https://ru.wikipedia.org/wiki/%D0%90%D0%92%D0%9B-%D0%B4%D0%B5%D1%80%D0%B5%D0%B2%D0%BE?ysclid=mnbu2k1jy3909571751}
\item \url{https://habr.com/ru/articles/150732/}
\end{itemize} 

\subsection*{Описание программы}

Программа реализована в 1 файле main.cpp

Описание основных типов данных:
\begin{enumerate}
\item \texttt{struct Node} — узел AVL-дерева, содержащий ключ (char*), значение (uint64\_t), высоту, указатели на левого и правого потомка.
\item \texttt{class Dictionary} — класс, инкапсулирующий AVL-дерево. Содержит корневой узел и методы для вставки, удаления, поиска, сохранения и загрузки.
\item \texttt{enum ErrorCode} — коды ошибок операций (ERR\_OK, ERR\_OUT\_OF\_MEMORY, ERR\_FILE\_OPEN, ERR\_FILE\_WRITE, ERR\_FILE\_READ, ERR\_INVALID\_FORMAT).
\item \texttt{enum InsertResult} — результаты вставки (INSERT\_OK, INSERT\_EXIST, INSERT\_ERROR).
\end{enumerate}

Описание основных функций и методов:
\begin{enumerate}
\item \texttt{Node::Node} — конструктор узла, выделяющий память под ключ и инициализирующий поля.
\item \texttt{Node::~Node} — деструктор, освобождающий память ключа.
\item \texttt{Dictionary::insert} — публичный метод вставки ключа и значения в дерево.
\item \texttt{Dictionary::remove} — метод удаления узла по ключу с балансировкой.
\item \texttt{Dictionary::search} — метод поиска значения по ключу.
\item \texttt{Dictionary::save/load} — методы сохранения и загрузки дерева в бинарный файл.
\item \texttt{leftRotate/rightRotate} — приватные методы выполнения поворотов для балансировки.
\item \texttt{balanceNode} — приватный метод балансировки узла после вставки/удаления.
\item \texttt{main} — основная функция, читающая команды из stdin, вызывающая методы словаря и выводящая результат.
\end{enumerate}

\subsection*{Дневник отладки}

\begin{enumerate}

\item При вставке дублирующегося ключа дерево ломалось.

Решение: добавлена проверка на существование ключа, возврат INSERT\_EXIST.

\item Неправильный порядок поворотов при LR и RL случаях.

Решение: сначала малый поворот потомка, затем большой поворот корня.

\item Утечка памяти при удалении узла с двумя потомками.

Решение: копирование данных из минимального узла правого поддерева с последующим удалением.

\item Поиск возвращал неверное значение для ключей в нижнем регистре.

Решение: приведение всех ключей к нижнему регистру через toLower() перед операцией.

\item Файл не читался после сохранения из-за неверного magic-числа.

Решение: проверка magic = 0x41564C31 при загрузке для валидации формата.

\item segfault при удалении из пустого дерева.

Решение: проверка на nullptr в начале функции remove().

\item Некорректная высота узла после удаления приводила к ошибке балансировки.

Решение: вызов updateHeight() перед вычислением баланса в balanceNode().

\item Ошибка чтения файла при нечётном формате (повреждённые данные).

Решение: очистка временного дерева и возврат ERR\_FILE\_READ при ошибке.

\item Двойное освобождение памяти при загрузке файла: ключ удалялся в цикле чтения и затем в деструкторе.

Решение: удаление key через delete[] только при ошибке чтения, иначе память передаётся узлу дерева.

\item Переполнение стека при вставке отсортированных данных: рекурсия insert() уходила на глубину O(n).

Решение: корректная балансировка гарантирует высоту O(log n), но потребовалась проверка updateHeight() в каждом узле пути.

\item Чтение ключа длиной 0 байт приводило к выделению нулевого массива и неопределённому поведению.

Решение: добавлена проверка len > 0 \&\& len <= 256 перед выделением памяти под ключ.

\end{enumerate}


\subsection*{Недочёты}

\begin{enumerate}
\item При загрузке повреждённого файла дерево может остаться в неконсистентном состоянии.

Почему не исправлено: реализована базовая валидация (magic-число, длина ключа), полная защита от всех видов повреждений избыточна.

\item Функция \texttt{search} имеет инвертированную логику перехода по дереву: при \texttt{cmp < 0} идёт переход вправо вместо влево.
Это работает только потому, что поиск проверяет оба поддерева рекурсивно.

Почему не исправлено: функциональность не нарушена.

\item Отсутствует проверка на переполнение буфера \texttt{path} (512 байт) при чтении пути к файлу.

Почему не исправлено: по ТЗ предполагается, что пути к файлам разумной длины.


\end{enumerate}


\subsection*{Выводы}

AVL-дерево эффективно для реализации словаря с операциями вставки, удаления и поиска за $O(\log n)$.
Типовые задачи: реализация ассоциативных контейнеров, индексация данных, хранение отсортированных коллекций.

Сложность программирования — выше средней. Алгоритм требует понимания балансировки, четырёх случаев поворотов и корректной работы с памятью.

Возникшие проблемы: утечки памяти при удалении узлов, некорректная балансировка LR/RL случаев, дублирование ключей,
валидация формата файла, приведение регистра символов. Все проблемы успешно устранены, программа соответствует требованиям ТЗ.

\end{document}